
\documentclass{article}

\usepackage[preprint]{neurips_2026}

\usepackage[utf8]{inputenc}
\usepackage[T1]{fontenc}
\usepackage{hyperref}
\usepackage{url}
\usepackage{booktabs}
\usepackage{amsfonts}
\usepackage{amsmath}
\usepackage{amssymb}
\usepackage{nicefrac}
\usepackage{microtype}
\usepackage{xcolor}
\usepackage{graphicx}
\usepackage{array}
\usepackage{caption}
\usepackage{subcaption}
\usepackage{float}

\providecommand{\tightlist}{%
  \setlength{\itemsep}{0pt}\setlength{\parskip}{0pt}}

\title{Mixtures of Algorithmic Experts:\\
  Exact Mathematical Modules in Differentiable Networks}

\author{%
  Stuart Whipp\\
  Independent Research\\
  \texttt{swhipp87@gmail.com}
}

\begin{document}

\maketitle

\begin{abstract}
We study a class of \emph{algorithmic neural modules}: parameter-free PyTorch modules
that expose exact, bounded computations---primality detection, Pascal's triangle,
binomial expansion---through differentiable interfaces, as frozen fixed-weight networks
composable with learned components.
Rather than asking a supervised model to approximate these tasks from finite samples,
we encode the corresponding finite algorithm directly as frozen network structure.
On tasks where standard MLPs trained from finite data fail to recover exact discrete
structure, the corresponding algorithmic module is provably exact by construction over
its bounded integer domain.
We further demonstrate a mixture-of-experts (MoE) architecture in which a learned
router discovers, without direct expert supervision, that it should delegate to the
structured expert: when both experts output probabilities and are blended by the router,
the MoE inherits the sieve's exactness---achieving 100\% accuracy across all evaluation
ranges versus 81\% for a comparable blind MLP.
We additionally present a multi-expert setting in which a single router simultaneously
learns to allocate primality queries to a \texttt{DifferentiableSieve} expert and
polynomial queries to a \texttt{BinomialExpansion} expert, producing a task-specific
routing confusion matrix without query-type labels.
A third expert slot---differentiable forward chaining over Horn-clause knowledge
bases---is already implemented in the accompanying library (\texttt{polyweave.reasoning})
and constitutes a natural next step toward neural systems that route between arithmetic,
algebraic, and logical reasoning.
All modules are released as \texttt{polyweave.maths}
(\url{https://github.com/swhipp/autodidact-hypernet}).
\end{abstract}

% ─────────────────────────────────────────────────────────────────────────────
\section{Introduction}
\label{sec:intro}
% ─────────────────────────────────────────────────────────────────────────────

The Peano axioms establish the natural numbers from two primitives: zero and a successor
function. Every arithmetic fact about $\mathbb{N}$ follows by induction from these
primitives alone. Our observation, and the starting point of this paper, is that the
same incremental construction is available inside a neural network. A single frozen
value and a fixed local rule---encoded as a weight matrix---suffice to generate Pascal's
triangle, Bernoulli's triangle, and the binomial coefficients of any $(Ax + By)^n$,
exactly and without training.

This matters for a practical reason. There is growing interest in equipping language
models with the ability to call external tools---calculators, symbolic solvers, code
interpreters---for tasks that require exact computation \citep{toolformer2023,
react2022}. We argue that a complementary capability is equally important: the ability
to perform a \emph{finite, well-defined class} of arithmetic operations \emph{internally},
before incurring the latency and brittleness of an external call. This is analogous
to human mental arithmetic---a person does not reach for a calculator to square a small
integer---and we suggest that structured neural modules provide a substrate for this kind
of reliable internal computation in artificial systems.

A natural question is: why not simply compute these results symbolically and look them
up? The answer is composability. A frozen sieve buffer inside a neural network is a
differentiable feature extractor whose output flows into any downstream learned
component---a classifier, a router, a language-model head. The structured module handles
the exact part; the learned component handles variation, uncertainty, and open-ended
generalisation. This asymmetry is the contribution: not that the algorithm can be
computed, but that it can be placed as a frozen differentiable expert inside a neural
routing system and discovered by a learned router without expert-label supervision.

Our modules are implemented as parameter-free \texttt{nn.Module} subclasses in PyTorch.
They expose a differentiable interface and can be composed with learned components
inside any gradient-based training loop.
A brief note on what we are and are not claiming: the point is not that these
algorithms \emph{cannot} be approximated by a large enough MLP with a better encoding.
The point is subtler and stronger---a finite sample of examples provides no guarantee
of exact extrapolation, whereas the structured module is exact within its bounded
integer domain by construction.
We contribute:

\begin{enumerate}
  \item Three zero-parameter modules---\texttt{PascalTriangle},
    \texttt{BinomialExpansion}, and \texttt{DifferentiableSieve}---that are provably
    exact over their bounded integer domain with no training data
    (Section~\ref{sec:modules}).
  \item A \texttt{BernoulliTriangle} module that encodes simplex and cake-number
    sequences as a column-wise cumulative sum of Pascal, illustrating the
    combinatorial richness accessible from a single fixed-weight primitive
    (Section~\ref{sec:modules}).
  \item A two-expert MoE in which a learned router discovers, without expert-label
    supervision, that it should delegate to the structured sieve; when both experts
    output probabilities the MoE achieves 100\% accuracy across all input ranges
    versus 81\% for a comparable blind MLP (Section~\ref{sec:tier2}).
  \item A multi-expert MoE routing simultaneously between a primality expert and a
    binomial-expansion expert, producing a task-specific routing confusion matrix
    learned purely from task losses (Section~\ref{sec:tier2b}).
  \item Release of all modules as \texttt{polyweave.maths}. The existing
    \texttt{polyweave.reasoning} forward chainer---differentiable Horn-clause inference
    via product t-norm---constitutes a concrete third expert type for future extension
    to logical reasoning (Section~\ref{sec:discussion}).
\end{enumerate}

\begin{figure}[t]
  \centering
  \includegraphics[width=0.82\linewidth]{plots/paper3_architecture.pdf}
  \caption{Architecture of the \texttt{polyweave.maths} module pool and the Tier-2
    mixture-of-experts system. The \texttt{DifferentiableSieve} (solid border) is the
    active structured expert; \texttt{BinomialExpansion} (dashed) is an additional
    passive module in the pool. A learned router selects between the sieve expert and a
    blind MLP. The \texttt{ForwardChainer} (dashed, right) represents the natural future
    third expert for logical inference, already implemented in
    \texttt{polyweave.reasoning}.}
  \label{fig:architecture}
\end{figure}

% ─────────────────────────────────────────────────────────────────────────────
\section{Related Work}
\label{sec:related}
% ─────────────────────────────────────────────────────────────────────────────

\paragraph{Neural arithmetic.}
Neural Arithmetic Logic Units \citep{nalu2018} learn to perform arithmetic operations
through weight-constrained modules. Neural GPU \citep{neuralgpu2015} solves
multiplication via convolutional recurrence. Our approach is complementary: rather than
learning to approximate arithmetic, we encode exact algorithms as frozen network
structure.

\paragraph{Differentiable programming and neurosymbolic reasoning.}
DeepProbLog \citep{deepproblog2018}, Logical Neural Networks \citep{lnn2020}, and
Neural Theorem Provers \citep{ntp2017} embed symbolic reasoning in differentiable
systems. \citet{differentiableILP2017} learn inductive logic programs end-to-end. We
focus on the numerical and combinatorial side of this space rather than logical
inference, though our accompanying \texttt{polyweave.reasoning} module bridges toward
the latter.

\paragraph{Mixture of experts.}
Standard MoE architectures \citep{moe2017,switch2022} route between learned experts.
We route between a \emph{frozen structured expert} and a learned expert---an asymmetric
setting in which the router must discover that exact structure exists, not merely which
of two approximators performs better.

\paragraph{Tool use and internal computation.}
Toolformer \citep{toolformer2023} and ReAct \citep{react2022} teach language models to
invoke external tools for arithmetic and search. Our work proposes an orthogonal
direction: certain computations should be handled \emph{internally} by structured
modules, reserving external calls for problems outside their scope.

\paragraph{This paper in context.}
This is the third in a series. Paper~1 \citep{paper1} established that multiplicative
(Sigma-Pi) neurons are recruited in proportion to the nonlinearity of the mapping they
are asked to learn, and introduced \texttt{pi\_scale} as a diagnostic. Paper~2
\citep{paper2} showed that transformer feed-forward blocks are largely linearly
recoverable and can be compressed into a single polynomial layer. The present paper
extends the \emph{structured-computation} theme to exact arithmetic modules.

% ─────────────────────────────────────────────────────────────────────────────
\section{Mathematical Modules}
\label{sec:modules}
% ─────────────────────────────────────────────────────────────────────────────

All modules are parameter-free \texttt{nn.Module} subclasses. Weights are registered
as frozen buffers; no gradient flows through them. Each module can be placed inside a
larger differentiable graph and is compatible with \texttt{torch.compile}.

\subsection{Pascal's Triangle (\texttt{PascalTriangle})}

Pascal's triangle satisfies $C(n,k) = C(n{-}1,k{-}1) + C(n{-}1,k)$.
We implement this as a fixed-weight 2-D convolution with kernel
$\mathbf{K} = \bigl[\begin{smallmatrix}1&1\\0&0\end{smallmatrix}\bigr]$,
applied iteratively with top-left zero-padding. Summing all intermediate states gives
the full triangle because intermediate $\mathbf{x}_k$ is non-zero only at row $k$:

\begin{equation}
  \mathbf{x}_{k+1}[i,j] = \mathbf{x}_k[i{-}1,j{-}1] + \mathbf{x}_k[i{-}1,j],
  \quad
  \text{output} = \sum_{k=0}^{N-1} \mathbf{x}_k.
\end{equation}

The output is exact (integer-valued floating-point) for $n \leq 15$.
Figure~\ref{fig:pascal} shows the result as a log-scale heatmap.

\begin{figure}[h]
  \centering
  \includegraphics[width=0.62\linewidth]{plots/paper3_pascal_heatmap.pdf}
  \caption{Pascal's triangle produced by \texttt{PascalTriangle(16)} — log-scale
    heatmap. Values annotated for $n \leq 7$. Zero learnable parameters; computed by
    iterative application of the fixed $2{\times}2$ convolution kernel
    $\mathbf{K} = \bigl[\begin{smallmatrix}1&1\\0&0\end{smallmatrix}\bigr]$.}
  \label{fig:pascal}
\end{figure}

\subsection{Binomial Expansion (\texttt{BinomialExpansion})}

Given scalars $A$, $B$ and exponent $n$, the binomial theorem gives:
\begin{equation}
  (Ax + By)^n = \sum_{k=0}^{n} \binom{n}{k} A^{n-k} B^k\, x^{n-k} y^k.
\end{equation}
We compute this by composing \texttt{PascalTriangle} with frozen exponent lookup tables.
The module takes $(A, B, n)$ and returns the exact coefficient vector of length
\texttt{num\_rows}, with trailing zeros beyond position $n$.
Zero learnable parameters; output matches symbolic computation to floating-point
precision on the supported integer domain.

\subsection{Bernoulli's Triangle (\texttt{BernoulliTriangle})}

Bernoulli's triangle contains the partial row sums of Pascal:
$B(n,k) = \sum_{j=0}^{k} C(n,j)$.
Notable sequences emerge along columns: naturals ($k{=}1$), 2-D lazy-caterer numbers
($k{=}2$: $1,2,4,7,11,\ldots$), 3-D cake numbers ($k{=}3$: $1,2,4,8,15,26,\ldots$),
and $k$-simplex numbers in general.
We implement $B$ as \texttt{torch.cumsum} applied column-wise to the Pascal output---a
natural fixed-weight operation consistent with the frozen-buffer pattern throughout.

\subsection{Differentiable Sieve (\texttt{DifferentiableSieve})}
\label{sec:sieve}

The Sieve of Eratosthenes marks composites by sweeping multiples of each prime
$p \leq \lfloor\sqrt{N}\rfloor$.
We implement each prime's sweep as a frozen \emph{comb buffer}---a tensor of length
$N{+}1$ with 1s at positions $0, p, 2p, \ldots$ (position $p$ cleared, since $p$ is
prime). Composite scores are combined via the probabilistic OR formula:
\begin{equation}
  \text{composite}(n) = 1 - \prod_{p \in \mathcal{P}} \bigl(1 - \text{comb}_p(n)\bigr),
\end{equation}
where $\mathcal{P}$ is the set of base primes. A final exponential decay
$s(n) = \exp(-\alpha \cdot \text{composite}(n))$ maps score 0 to 1 (prime) and 1 to 0
(composite). The decay parameter $\alpha$ controls sharpness and may be made learnable.
Detection is exact for all $n \leq N$ when $\max_p = \lfloor\sqrt{N}\rfloor$.

\subsection{Modular Membership via \texttt{radbas}}

A recurring primitive is soft equality testing via the radial-basis activation
$r(x) = \exp(-(\varepsilon x)^2)$, which peaks at $x{=}0$ \citep{paper1}.
For membership queries---``is $n$ a multiple of $p$?''---we precompute known multiples
as a lookup table, subtract the query, and apply \texttt{radbas}. The score is near 1
for exact matches and near 0 elsewhere, with no training required.

% ─────────────────────────────────────────────────────────────────────────────
\section{Experiments}
\label{sec:experiments}
% ─────────────────────────────────────────────────────────────────────────────

\subsection{Tier-1: Zero-Shot Structural Tasks}
\label{sec:tier1}

We evaluate three tasks for which the structured module achieves provable 100\% accuracy
without any training data. In each case we compare against an MLP baseline trained with
Adam \citep{adam2015} on the same task.

\paragraph{Task A --- Primality detection.}
Input: integer $n \in [2, 500]$.
Label: \texttt{is\_prime}$(n)$.
The structured module is \texttt{DifferentiableSieve(500)}.
The MLP baseline (3 hidden layers, 128 units, ReLU) is trained on $n \in [2, 100]$ for
300 epochs and evaluated on the full range.
Results in Table~\ref{tab:tier1} and Figure~\ref{fig:sieve} show that the sieve
is provably exact over its bounded integer domain; the MLP trained on $n \in [2,100]$
reaches 76.8\% in-range and 82.5\% out-of-range---failing to recover exact primality
from finite samples, and requiring 7.97\,s of training and 33{,}409 parameters to do
so. The MLP baseline is illustrative rather than definitive: a larger model with a
richer input encoding might do better. The stronger claim is that \emph{no sampled
approximator} provides a structural guarantee of exact extrapolation, whereas the
sieve does so by construction.

\paragraph{Task B --- Binomial expansion.}
Input: $(A, B, n)$ with $A, B \in \{-8,\ldots,8\}\setminus\{0\}$, $n \in \{2,\ldots,8\}$.
Label: coefficient vector of $(Ax + By)^n \in \mathbb{Z}^{16}$.
The dataset is a deterministic sweep of 1{,}792 samples (all valid triples), split
80/20 for the MLP baseline.
The structured module (\texttt{BinomialExpansion}) achieves 100\% exact accuracy on
all samples with zero training. The MLP (4 hidden layers, 256 units) achieves 0.8\%
exact accuracy on the validation split after 500 epochs: it approximates the
leading-coefficient magnitude but cannot recover integer polynomial coefficients
reliably. Figure~\ref{fig:binomial} shows the failure mode directly.

\begin{table}[t]
  \centering
  \caption{Tier-1 results. \emph{Exact} = fraction of outputs matching ground truth
    to integer precision. Sieve accuracy is provable; MLP accuracy is empirical.}
  \label{tab:tier1}
  \small
  \begin{tabular}{llrrrrr}
    \toprule
    Task & Method & Params & Train (s) & In-range & OOD & Exact (val) \\
    \midrule
    Primality   & \texttt{DifferentiableSieve} & \textbf{0} & \textbf{0} & \textbf{100.0\%} & \textbf{100.0\%} & --- \\
    Primality   & MLP baseline                 & 33{,}409   & 7.97       & 76.8\%            & 82.5\%            & --- \\
    \midrule
    Binomial    & \texttt{BinomialExpansion}   & \textbf{0} & \textbf{0} & ---               & ---               & \textbf{100.0\%} \\
    Binomial    & MLP baseline                 & 202{,}512  & 33.60      & ---               & ---               & 0.8\%  \\
    \bottomrule
  \end{tabular}
\end{table}

\begin{figure}[t]
  \centering
  \begin{subfigure}[b]{0.48\linewidth}
    \includegraphics[width=\linewidth]{plots/paper3_tier1_primality.pdf}
    \caption{Primality detection: sieve scores (top) vs MLP sigmoid outputs (bottom)
      for $n \in [2, 500]$. Red = prime, blue = composite. Dashed line = training
      boundary. The sieve separates cleanly everywhere; the MLP produces diffuse scores
      throughout.}
    \label{fig:sieve}
  \end{subfigure}
  \hfill
  \begin{subfigure}[b]{0.48\linewidth}
    \includegraphics[width=\linewidth]{plots/paper3_tier1_binomial.pdf}
    \caption{Binomial expansion: structured module outputs (exact) vs MLP predictions
      on the validation split. The MLP approximates magnitude but cannot recover
      integer coefficients; exact accuracy is 0.8\%.}
    \label{fig:binomial}
  \end{subfigure}
  \caption{Tier-1 zero-shot results. In both tasks the structured module achieves 100\%
    exact accuracy with zero parameters and zero training time.}
  \label{fig:tier1}
\end{figure}

\subsection{Tier-2: MoE Range Extrapolation}
\label{sec:tier2}

We train a two-expert primality MoE consisting of:
(i)~a \texttt{DifferentiableSieve} expert (0 learnable parameters; output
$\hat{p}_\text{sieve}(n) = \exp(-\alpha \cdot \text{composite}(n)) \in (0,1)$,
detached from the gradient graph);
(ii)~a blind MLP expert (2 hidden layers, 64 units; output converted to probability
via sigmoid before blending);
(iii)~a learned router MLP (1 hidden layer, 32 units; softmax output $\mathbf{w}$).
The combined output is the probability mixture $\hat{p} = w_0 \hat{p}_\text{sieve}
+ w_1 \hat{p}_\text{blind}$, trained with binary cross-entropy directly on $\hat{p}$.
Blending probabilities---rather than mixing a probability with a raw logit---is the
key design choice: it ensures the router can collapse to $w_0 = 1$ and thereby fully
inherit the sieve's accuracy.

Starting from near-uniform initialisation ($w_0 \approx 0.58$), the router weight
rises monotonically, settling at $w_0 = 0.996$ by epoch 400.
Table~\ref{tab:tier2} shows that the MoE achieves 100\% accuracy across both
in-range and out-of-distribution ranges---matching the sieve exactly rather than
averaging with an inferior MLP signal.

\begin{table}[t]
  \centering
  \caption{Tier-2 MoE: single-task primality. Probability blending allows the MoE
    to fully inherit sieve accuracy at $w_0 = 0.996$.}
  \label{tab:tier2}
  \small
  \begin{tabular}{lrrrr}
    \toprule
    Method & Params (learned) & In-range [2,100] & OOD [101,500] & All [2,500] \\
    \midrule
    \texttt{DifferentiableSieve}         & \textbf{0} & \textbf{100.0\%} & \textbf{100.0\%} & \textbf{100.0\%} \\
    MoE: sieve + MLP (prob mix)          & 4{,}483    & \textbf{100.0\%} & \textbf{100.0\%} & \textbf{100.0\%} \\
    Blind MLP                            & 4{,}353    & 74.8\%           & 82.5\%           & 81.0\%           \\
    \bottomrule
  \end{tabular}
\end{table}

\begin{figure}[t]
  \centering
  \includegraphics[width=\linewidth]{plots/paper3_tier2_moe.pdf}
  \caption{Tier-2 (single-task) MoE. \emph{Left}: router weight $w_0$ toward the
    sieve expert, rising from ${\approx}0.58$ to $0.996$ once both experts output
    probabilities. \emph{Centre}: per-$n$ router allocation at inference.
    \emph{Right}: accuracy by range.}
  \label{fig:router}
\end{figure}

\subsection{Tier-2b: Multi-Expert Routing}
\label{sec:tier2b}

We extend the MoE to two structured experts and two task types, asking whether a
single router can learn task-to-expert allocation from task losses alone---without
being told which expert handles which task.

\paragraph{Setup.}
The expert pool contains:
(0)~\texttt{DifferentiableSieve} (primality, frozen);
(1)~\texttt{BinomialExpansion} (polynomial coefficients, frozen);
(2)~\texttt{DualHeadMLP} (learned residual for both tasks).
A single router MLP takes a 6-dimensional query feature---encoding task type and
query parameters in its first two and remaining four dimensions respectively---and
outputs softmax weights $(w_0, w_1, w_2)$ over the three experts.
Training uses a joint loss: binary cross-entropy on primality outputs and MSE on
normalised coefficient vectors.

Crucially, the structured expert not relevant to a given task contributes nothing to
that task's output directly. The router is nonetheless penalised for allocating mass
to the irrelevant expert via the softmax Jacobian: increasing $w_\text{binom}$ on a
primality query reduces $w_\text{sieve}$, raising the primality loss.

\paragraph{Results.}
After 600 epochs, the router collapses almost entirely to the correct expert
(Table~\ref{tab:confusion}; Figure~\ref{fig:confusion}):

\begin{table}[h]
  \centering
  \caption{Multi-expert routing confusion matrix (mean gate weight per task).
    Learned from task losses only---no expert-identity labels provided.}
  \label{tab:confusion}
  \small
  \begin{tabular}{lrrr}
    \toprule
    Task & \texttt{DifferentiableSieve} & \texttt{BinomialExpansion} & MLP \\
    \midrule
    Primality ($n \in [2,100]$)                         & \textbf{1.000} & 0.000 & 0.000 \\
    Binomial ($A,B \in \{1..4\}$, $n \in \{2..5\}$)   & 0.056 & \textbf{0.931} & 0.013 \\
    \bottomrule
  \end{tabular}
\end{table}

Under top-1 (hard) routing---using the argmax expert's output directly---both tasks
achieve 100\% exact accuracy on training and OOD ranges, since the router selects the
exact structured expert in virtually all cases.
This demonstrates that a single learned router can discover and exploit multiple
structured experts simultaneously, partitioning by task type purely from gradient
signal.

\begin{figure}[t]
  \centering
  \includegraphics[width=\linewidth]{plots/paper3_tier2_multi_expert.pdf}
  \caption{Multi-expert routing. \emph{Left}: routing confusion matrix (green
    intensity = mean gate weight); the router allocates primality to
    \texttt{DifferentiableSieve} ($w_0 = 1.000$) and binomial queries to
    \texttt{BinomialExpansion} ($w_1 = 0.931$) without task-identity labels.
    \emph{Right}: router weight evolution per task type; within ${\sim}100$ epochs
    the correct task partition has been discovered.}
  \label{fig:confusion}
\end{figure}

% ─────────────────────────────────────────────────────────────────────────────
\section{Discussion}
\label{sec:discussion}
% ─────────────────────────────────────────────────────────────────────────────

\paragraph{Why not just compute it symbolically?}
The point is composability and differentiability. A frozen sieve buffer inside a neural
network serves as a feature extractor whose output flows into a learned classifier or
router. The structured module handles the exact part; the learned component handles
variation and uncertainty around it. This asymmetry mirrors how humans combine mental
arithmetic (fast, exact, bounded domain) with estimation or tool use (slower, flexible,
unbounded).

\paragraph{Relationship to the Peano axioms.}
Our modules demonstrate that combinatorial structures central to mathematics emerge from
Peano-like primitives: a single seed value and a fixed local rule, encoded as frozen
network weights. \texttt{PascalTriangle} starts from a single 1 and derives every
binomial coefficient by one step of the successor-like convolution
$\mathbf{K} = \bigl[\begin{smallmatrix}1&1\\0&0\end{smallmatrix}\bigr]$.
\texttt{BernoulliTriangle} applies a column-wise cumulative sum to the result.
\texttt{DifferentiableSieve} sweeps multiples of each prime in the same incremental
spirit. The underlying claim---that rich mathematical structure can be bootstrapped from
minimal local rules---is as old as Peano and as new as the observation that the same
operations are expressible as differentiable network layers.

\paragraph{Extension to logical inference.}
The \texttt{polyweave} library already ships a differentiable forward chainer over
propositional Horn knowledge bases (\texttt{polyweave.reasoning.ForwardChainer}), which
uses the product t-norm as a differentiable AND gate \citep{paper1}.
This module constitutes a natural third arm of the MoE: a \emph{logic expert} that
handles entailment queries exactly, alongside the arithmetic experts demonstrated here.
A router would then learn to distinguish numeric from logical queries---a step toward
the unified differentiable reasoning system sketched in Figure~\ref{fig:architecture}.
We leave this combination to future work, noting that all components are already
implemented and publicly available.

\paragraph{Toward a Large Maths Model.}
The MoE pattern demonstrated here scales to a richer expert pool: modular arithmetic
(\texttt{radbas} membership), recurrence relations (fixed-weight RNN), sorting
(differentiable sorting networks), and GCD (bounded Euclidean unrolling) are all
expressible as parameter-free modules in the same framework. Combined with a
natural-language front-end---analogous to the convolutional LaTeX parser prototyped
alongside this work---such a system would constitute a \emph{Large Maths Model}: an
architecture in which a language model delegates well-posed arithmetic queries to
exact internal modules rather than approximating them in its feed-forward layers.

\paragraph{Limitations.}
The modules are exact only within their specified ranges ($N$ for the sieve;
\texttt{num\_rows} for Pascal). The router weight in the two-expert MoE
(Section~\ref{sec:tier2b}) does not collapse entirely to $w_\text{binom} = 1.0$
for binomial queries ($w_1 = 0.931$) because the joint training loss balances both
tasks; extended training or a higher learning rate annealing schedule may close this
gap further. Extending the expert pool to continuous domains or general recurrence
relations remains future work.

% ─────────────────────────────────────────────────────────────────────────────
\section{Conclusion}
\label{sec:conclusion}
% ─────────────────────────────────────────────────────────────────────────────

We have shown that classical combinatorial algorithms are expressible as zero-parameter
neural network modules, enabling exact computation within differentiable architectures.
Learned routers, trained only on task losses, discover without supervision which
structured expert to delegate each task type to---including a multi-expert setting
where a single router simultaneously allocates primality queries to the sieve and
polynomial queries to the binomial module, achieving 100\% exact accuracy on both
in-distribution and out-of-distribution inputs under hard routing.
The three modules presented here, together with the existing \texttt{polyweave.reasoning}
forward chainer, form the nucleus of a broader differentiable reasoning library
whose capabilities extend from arithmetic through to propositional logic.
We release all code and experiments at \url{https://github.com/swhipp/autodidact-hypernet}.

% ─────────────────────────────────────────────────────────────────────────────
\section*{Acknowledgements}

The author thanks the open-source communities behind PyTorch, NumPy, and Matplotlib.

% ─────────────────────────────────────────────────────────────────────────────
\bibliographystyle{plainnat}
\bibliography{references_paper3}

% ─────────────────────────────────────────────────────────────────────────────
\appendix

\section{Module Implementation Details}
\label{app:impl}

All modules are available as \texttt{polyweave.maths.\{PascalTriangle, BinomialExpansion,
BernoulliTriangle, DifferentiableSieve\}}. Each inherits from \texttt{nn.Module} and
registers its structural weights as frozen buffers via \texttt{register\_buffer}.
The \texttt{forward()} method is decorated with \texttt{@torch.no\_grad()} where no
downstream gradient is required.

\paragraph{PascalTriangle.} $O(N^2)$ memory; $O(N)$ forward passes of a single
$1{\times}1{\times}2{\times}2$ convolution. Exact for $n \leq 15$ in \texttt{float32}.

\paragraph{DifferentiableSieve.} $O(N \log \log N)$ buffer storage (by the prime
harmonic series). Decay parameter $\alpha$ defaults to 5.0; values $\geq 10$ give
near-binary output.

\paragraph{BernoulliTriangle.} Single \texttt{torch.cumsum} call on the Pascal output;
negligible additional cost.

\section{Extended Results}
\label{app:results}

\begin{table}[h]
  \centering
  \caption{Tier-1 primality: per-range breakdown.}
  \small
  \begin{tabular}{lrrrr}
    \toprule
    Method & [2, 50] & [51, 100] & [101, 250] & [251, 500] \\
    \midrule
    \texttt{DifferentiableSieve} & 100.0\% & 100.0\% & 100.0\% & 100.0\% \\
    MLP baseline & [TBD] & [TBD] & [TBD] & [TBD] \\
    \bottomrule
  \end{tabular}
\end{table}

\end{document}
